\documentclass[10pt]{article}

\def\Type{LessonDJ}
\def\TypeNum{Workshop Week 13} %Lesson #
\def\TypeTitle{FTOC and Intro to Probability} %Lesson Title
\def\Semester{Fall 2021} %Not Used 
\def\Course{MATH 1190}%Course number and name
\def\Targets{    \target{I2}, \target{I3}, \target{F3} }

\usepackage{preambleDJ} 

%%%%%%%%%%%% BEGIN DOCUMENT %%%%%%%%%%%%%%%
%%%%%%%%%%%% BEGIN DOCUMENT %%%%%%%%%%%%%%%
%%%%%%%%%%%% BEGIN DOCUMENT %%%%%%%%%%%%%%%
%%%%%%%%%%%% BEGIN DOCUMENT %%%%%%%%%%%%%%%
%%%%%%%%%%%% BEGIN DOCUMENT %%%%%%%%%%%%%%%

\begin{document}

A manufacturing company needs to design an open-top rectangular box with a square base. The box must have a fixed volume of $4000 \text{ cm}^3$. The material for the base costs $5$ cents per $ \text{cm}^2$, and the material for the sides costs $2$ cents per $ \text{cm}^2$. Find the dimensions of the box that minimize the total cost of materials.

\vspace{2mm}
\textbf{Modeling Phase:}
\begin{enumerate}
    \item[(a)] Let $x$ be the length of the sides of the square base, and $h$ be the height of the box. Write down the formula for the area of the base $A_{\text{base}}$ and the formula for the area of the sides $A_{\text{sides}}$.
    
    \item[(b)] Write down the formula for the total cost $C$ of the box in terms of $x$ and $h$, given that
    $$
        C = (\text{The base cost per cm}^2) \times A_{\text{base}} + (\text{The sides cost per cm}^2) \times A_{\text{sides}}
    $$
    \item[(c)] Write down the constraint equation based on the required volume of the box.
    \item[(d)] Use the volume constraint to express $h$ in terms of $x$, then substitute your expression for $h$ into the cost equation to obtain a single-variable function $C(x)$.
\end{enumerate}

\vfill
\textbf{Calculus Phase:}
Now that you have obtained $C(x)$ as a single variable function, use calculus to find the absolute minimum cost and the corresponding dimensions $x$ and $h$.
\vfill
\newpage

You are the manager of an apartment complex with $50$ units. When you set the rent at $\$800$ per month, all apartments are rented. As you increase the rent by $\$25$ per month, one fewer apartment is rented. Maintenance costs run $\$50$ per month for each occupied unit. What is the rent that maximizes the total amount of profit?

\vspace{2mm}
\textbf{Modeling Phase:}
\begin{enumerate}
    \item[(a)] Let $x$ represent the number of increases by $\$25$. Write down the formula for the rent per unit $R(x)$ and the number of occupied units $N(x)$.
    \item[(b)] Write down the total revenue in terms of $x$, given that
    $$
        \text{Revenue} = (\text{Rent per unit}) \times (\text{Number of occupied units}).
    $$
    \item[(c)] Write down the formula for the total maintenance cost.
    \item[(d)] Write down the formula for the profit, given that
    $$
        \text{Profit} = \text{Revenue} - \text{Cost}.
    $$
\end{enumerate}
\vfill

\textbf{Calculus Phase:}
Now that you have obtained the formula for the profit as a single variable function, use calculus to find the rent that maximizes the total amount of profit.
\vfill

\end{document}